%% kapa: exam=1 round=32 year=2021 period=1 form=A subject=economics subject_ko=경제학원론 seq=24 number=64
%% source: Q-Net 감정평가사 기출문제 (cst00309) artlSeq=5208397 file=extracted/1008/5208397/A형/2021년도 제32회 감정평가사 1차 1교시 A형.pdf
%% conversion: kordoc 4.24.51 -> tools/gen1.py (auto); review=proofread-sonnet; flags=none
\begin{question}{64}
폐쇄경제에서 국내총생산이 소비, 투자, 그리고 정부지출의 합으로 정의된 항등식이 성립할 때, 국내총생산과 대부자금시장에 관한 설명으로 옳지 않은 것은?
\begin{choices}
  \item 총저축은 투자와 같다.
  \item 민간저축이 증가하면 투자가 증가한다.
  \item 총저축은 민간저축과 정부저축의 합이다.
  \item 민간저축이 증가하면 이자율이 하락하여 정부저축이 증가한다.
  \item 정부저축이 감소하면 대부시장에서 이자율은 상승한다.
\end{choices}
\end{question}
\answer{4}
